\documentclass[11pt,a4paper]{article}
\usepackage[T1]{fontenc}
\usepackage{lmodern}
\usepackage[margin=25mm]{geometry}
\usepackage{microtype,amsmath,amssymb,amsthm,mathtools}
\usepackage{booktabs,longtable,array,enumitem,xcolor}
\usepackage{xurl}
\usepackage{hyperref}
\usepackage{fancyhdr}
\definecolor{ink}{RGB}{28,49,72}
\hypersetup{colorlinks=true,linkcolor=ink,citecolor=ink,urlcolor=ink,pdftitle={Two Fourier Peaks and a Moving Lattice Maximum},pdfauthor={Research report prepared for Vladimir Reshetnikov}}
\pagestyle{fancy}\fancyhf{}\fancyhead[L]{\small Two Fourier peaks and a lattice maximum}\fancyhead[R]{\small OEIS A039831}\fancyfoot[C]{\thepage}
\setlength{\headheight}{14pt}
\setlength{\emergencystretch}{2em}
\numberwithin{equation}{section}
\newtheorem{theorem}{Theorem}[section]
\newtheorem{lemma}[theorem]{Lemma}
\newtheorem{proposition}[theorem]{Proposition}
\newtheorem{corollary}[theorem]{Corollary}
\theoremstyle{definition}\newtheorem{definition}[theorem]{Definition}
\theoremstyle{remark}\newtheorem{remark}[theorem]{Remark}
\newcommand{\E}{\mathbb E}\newcommand{\PP}{\mathbb P}
\newcommand{\Z}{\mathbb Z}\newcommand{\R}{\mathbb R}
\newcommand{\He}{\operatorname{He}}
\newcommand{\dist}{\operatorname{dist}}
\newcommand{\ee}{\mathrm e}
\newcommand{\ii}{\mathrm i}
\newcommand{\dd}{\mathrm d}
\newcommand{\calE}{\mathcal E}
\newcommand{\A}{\mathcal A}
\newcommand{\B}{\mathcal B}
\newcommand{\HH}[2]{H_{#1}^{(#2)}}
\title{\textbf{Two Fourier Peaks and a Moving Lattice Maximum}\\[5pt]
\large A proof of the OEIS A039831 asymptotic conjecture,\\
four correction orders, a smoothing hierarchy, and inverse growth}
\author{Research report prepared for Vladimir Reshetnikov}
\date{October 1, 2026}
\begin{document}
\maketitle
\begin{abstract}
Let $M_n$ be the largest coefficient of
$P_n(q)=\prod_{j=1}^n(1+q+q^3+\cdots+q^{2j-1})$.
The OEIS entry A039831 records V\'aclav Kote\v{s}ovec's 2023 conjecture
$M_n\sim3(n/\ee)^n$. We prove this statement and obtain explicit corrections
through relative order $n^{-4}$, together with an algorithmic expansion to
arbitrary algebraic order. Two features distinguish the answer from a routine
Gaussian estimate. Integer rounding produces a nonconstant log-periodic term
at order $n^{-3}$. A second Fourier neighborhood, at $q=-1$, produces a
parity term at order $n^{-4}$, despite the exact equality of the total even-
and odd-degree coefficient masses. This smaller contribution changes the
maximizing exponent infinitely often; the number of failures of a
single-neighborhood rounding rule up to $N$ has order $\log N$.
We prove a zero-multiplicity law for suppression of the second neighborhood
under fixed Bernoulli smoothing, and derive Lambert-$W$ inverse-growth
formulas with the integer-threshold qualification made explicit.
Exact coefficient computations through $n=400$ and symbolic coefficient
checks accompany the article. The arguments are self-contained mathematical
proofs, not Lean-checked declarations. The source audit establishes the
conjecture's current OEIS wording, not an exhaustive claim of publication
priority.
\end{abstract}

\paragraph{Editorial note (October 1, 2026, batch 73O2).}
This report is manuscript~62 of batch~73 of ProveIt's incoming-reports
intake, placed in the research-report collection in commit
\texttt{6e193dd4f}. It is AI-assisted, unrefereed and not formalized. The
text is the delivered manuscript; the write added the label prefix
\texttt{tfp:}, notes marked ``[Added October 1, 2026, batch 73O2]'', one
bibliography entry, and Appendix~\ref{tfp:app:provenance}. The distinct-coefficient
conjecture of A039824 is \emph{not} resolved here.
\newpage
\tableofcontents
\newpage

\section{Scope, provenance, and the main result}
\subsection{The problem and why the higher orders matter}
Write
\begin{equation}\label{tfp:eq:P}
 P_n(q)=\prod_{j=1}^n\left(1+\sum_{a=1}^j q^{2a-1}\right)
       =\sum_{k=0}^{n^2}c_{n,k}q^k,
 \qquad M_n=\max_{k\in\Z}c_{n,k},
\end{equation}
with coefficients outside $[0,n^2]$ taken to be zero. This is the definition
of OEIS A039831. Its record, inspected on October 1, 2026, still labels
$M_n\sim3n^n/\ee^n$ a conjecture and attributes it to V\'aclav Kote\v{s}ovec,
January 5, 2023 \cite{oeis831}. The related entry A039824 concerns the
\emph{number of distinct coefficient values}, a different question; its
conjectured eventual formula $n^2-3$ is not proved here \cite{oeis824}.

The product in \eqref{tfp:eq:P} is a simple, nonnegative, growing-support
convolution model. Its total mass is factorial, its width is of order
$n^{3/2}$, and its center has a harmonic correction. It is useful as an
explicit test case for a general issue: a coefficient maximum is sensitive
both to the integer lattice and to small Fourier contributions that a
central limit theorem does not resolve. The leading conjecture alone is an
application of local-limit methods. The substantive refinements below are
the explicit lattice-resolved expansion, its mode transitions, and the
suppression law under smoothing.

\subsection{Relationship to ProveIt and limits of the novelty claim}
The inspected ProveIt materials distinguish research arguments, executable
checks, and the formal trusted theorem boundary \cite{proveit}. This report
follows that distinction. It connects to the repository's transseries and
inversion program, particularly its separation of a logarithmic core,
remainder transport, and staircase inversion \cite{transseries}. A targeted
repository search for A039831 returned no match. This is not a claim that
every file in the repository was examined.

The Fourier method itself is classical. Modern work of Dolgopyat and
Hafouta explicitly treats trigonometric corrections and the instability of
modular uniformity after deletion of finitely many summands
\cite{dolgopyat}. Their stated setting uses uniformly bounded summands;
our summands have growing support. We do not invoke their theorem outside
its hypotheses. Instead we prove the required bounds directly and use
that work to situate the mechanism. Targeted A-number searches did not
identify an earlier proof of the particular conjecture or the refinements
proved here. These searches are not a comprehensive literature review,
and no global priority or ``major breakthrough'' claim is required for
the mathematical statements.

\subsection{Notation and a first theorem}
Let
\begin{equation}\label{tfp:eq:notation}
 h_n=H_{n+1}=\sum_{j=1}^{n+1}\frac1j,
 \quad h_{2,n}=\HH{n+1}{2}=\sum_{j=1}^{n+1}\frac1{j^2},
 \quad \mu_n=\frac{n(n-1)}2+h_n-1,
 \quad \calE_n=3(n/\ee)^n.
\end{equation}
The symbols $O(\cdot)$ below refer to $n\to\infty$; constants can depend
on a fixed expansion order or a fixed bounded range of $k-\mu_n$.
Define
\begin{align}
 a_1&=-\frac{431}{300}, &
 a_2&=\frac{23085971}{1764000},\label{tfp:eq:a12}\\
 C_3(d;h,h_2)&=-\frac92d^2-\frac{27}2h+\frac92h_2
                  -\frac{107751433393}{1587600000},\label{tfp:eq:C3}\\
 C_4(d;h,h_2)&=\frac{1263}{40}d^2-\frac{81}2d
       +\frac{3789}{40}h-\frac{1263}{40}h_2
       +\frac{155775441709528423}{322727328000000}.
       \label{tfp:eq:C4}
\end{align}

\begin{theorem}[Four correction orders and the conjecture]\label{tfp:thm:main}
For every fixed $D>0$, uniformly for integers $k$ satisfying
$|d|=|k-\mu_n|\le D$,
\begin{equation}\label{tfp:eq:fourorders}
\begin{split}
 \frac{c_{n,k}}{\calE_n}
  ={}&1+\frac{a_1}{n}+\frac{a_2}{n^2}
     +\frac{C_3(d;h_n,h_{2,n})}{n^3}\\
    &+\frac{C_4(d;h_n,h_{2,n})-18(-1)^{n+k}}{n^4}
     +O_D\!\left(\frac{1+\log^2 n}{n^5}\right).
\end{split}
\end{equation}
For all sufficiently large $n$, every maximizing exponent belongs to
$\{\lfloor\mu_n\rfloor,\lceil\mu_n\rceil\}$. Consequently
\begin{equation}\label{tfp:eq:leading}
 M_n=3(n/\ee)^n\left(1-\frac{431}{300n}
            +\frac{23085971}{1764000n^2}
            +O\!\left(\frac{\log n}{n^3}\right)\right).
\end{equation}
In particular, Kote\v{s}ovec's stated asymptotic conjecture is true.
The maximum of the two expressions in \eqref{tfp:eq:fourorders}, evaluated at
the indicated candidate exponents, approximates $M_n/\calE_n$ with the
same stated error, including at a mode transition.
\end{theorem}

The proof occupies Sections \ref{tfp:sec:prob}--\ref{tfp:sec:modes}. Exact harmonic
numbers in \eqref{tfp:eq:fourorders} are intentional: replacing them by their
limits without expanding their tails changes later coefficients. The
remainder is an asymptotic bound, not a numerical interval certificate
with an explicit threshold.

\section{Independent summands and exact cumulants}\label{tfp:sec:prob}
\subsection{Probability normalization}
Let $X_j$ be independent, each uniform on the $j+1$ integers
\begin{equation}
 \{0,1,3,\ldots,2j-1\},\qquad S_n=\sum_{j=1}^nX_j.
\end{equation}
Evaluation of \eqref{tfp:eq:P} at $q=1$ gives
\begin{equation}\label{tfp:eq:normalization}
 P_n(1)=(n+1)!,\qquad
 c_{n,k}=(n+1)!\PP(S_n=k).
\end{equation}
The moment generating function of a single summand can be written
\begin{equation}\label{tfp:eq:Fj}
 F_j(t)=1+\sum_{a=1}^j e^{(2a-1)t}
       =1+e^{jt}\frac{\sinh(jt)}{\sinh t},
 \qquad \E e^{tX_j}=\frac{F_j(t)}{j+1},
\end{equation}
where the quotient is interpreted by continuity at zero.

Let $\kappa_{r,n}$ denote the $r$th cumulant of $S_n$, and put
$V_n=\kappa_{2,n}$. Direct summation yields
\begin{align}
 \kappa_{1,n}&=\mu_n,\label{tfp:eq:mean}\\
 V_n&=\frac{n^3}{9}+\frac{n^2}{2}-\frac{29n}{18}
       +3h_n-h_{2,n}-2,\label{tfp:eq:variance}\\
 \kappa_{3,n}&=n^2-6n+14h_n-9h_{2,n}
          +2\HH{n+1}{3}-7.\label{tfp:eq:kappa3}
\end{align}
In particular $V_n\sim n^3/9$. This, together with factorial total mass,
already explains the proposed factor $3$.

\subsection{An exact cumulant algorithm}
For $r\ge1$, set
\begin{equation}\label{tfp:eq:moment}
 m_r(j)=\frac1{j+1}\sum_{a=1}^j(2a-1)^r,
 \qquad m_0(j)=1.
\end{equation}
The single-summand cumulants are generated by
\begin{equation}\label{tfp:eq:cumrec}
 b_r(j)=m_r(j)-\sum_{\ell=1}^{r-1}
       \binom{r-1}{\ell-1}b_\ell(j)m_{r-\ell}(j),
 \qquad \kappa_{r,n}=\sum_{j=1}^n b_r(j).
\end{equation}
This follows by differentiating $M'(t)=K'(t)M(t)$ for $M=e^K$ at zero.
For example,
\begin{align}
 b_2(j)&=\frac{j^2}{3}+\frac{2j}{3}-2
                  +\frac3{j+1}-\frac1{(j+1)^2},\\
 b_3(j)&=2j-7+\frac{14}{j+1}
                  -\frac9{(j+1)^2}+\frac2{(j+1)^3},\\
 b_4(j)&=-\frac{2j^4}{15}-\frac{8j^3}{15}-\frac{4j^2}{5}
       +\frac{172j}{15}-40+\frac{93}{j+1}
       -\frac{83}{(j+1)^2}+\frac{36}{(j+1)^3}
       -\frac6{(j+1)^4}.
\end{align}
Thus all fixed-order cumulants are explicitly expressible using
polynomials in $n$ and generalized harmonic numbers.

\begin{lemma}[Cumulant size and odd-order cancellation]\label{tfp:lem:cumulants}
For each fixed $r\ge2$, $\kappa_{r,n}=O_r(n^{r+1})$.
For odd $r\ge3$, the stronger estimate
$\kappa_{r,n}=O_r(n^{r-1})$ holds. The analogous estimates hold for the
second-neighborhood cumulants defined in Section \ref{tfp:sec:second}.
\end{lemma}
\begin{proof}
The first estimate follows either from \eqref{tfp:eq:cumrec} and
$|X_j|\le2j$, or by applying Cauchy estimates to the logarithm of
\eqref{tfp:eq:Fj} on a disk of radius a sufficiently small constant divided
by $j$.
For the cancellation, set $t=z/j$ and
$E(z)=z\coth z$, an even analytic function. Uniformly for small $z$,
\begin{equation}
 \frac{F_j(z/j)}{j+1}
 =\frac{e^z\sinh z/z}{1+1/j}
  \left(1+\frac{E(z)-z}{j}-\frac{z^2}{6j^2}
                +O(j^{-4})\right).
\end{equation}
The zeroth-order logarithm has only the linear odd term $z$.
The $j^{-1}$ logarithmic term has only the linear odd term $-z/j$.
All odd Taylor coefficients of degree at least three therefore first
occur at order $j^{-2}$. Taking $r$ derivatives with respect to $t$
shows $b_r(j)=O_r(j^{r-2})$, and summation gives the assertion.
For the signed factors used below, replace $E(z)-z$ by its negative and
$1+1/j$ by $1-1/j$; the same cancellation applies. The extra fixed
factor $(e^t-1)/t$ has bounded cumulants and causes no difficulty.
\end{proof}

\section{Why two Fourier neighborhoods are necessary}\label{tfp:sec:second}
\subsection{The neighborhood of \texorpdfstring{$q=-1$}{q = -1}}
Define an analytic function with value one at the origin by
\begin{equation}\label{tfp:eq:Bdef}
 B_n(t)=\frac{e^t-1}{t}
       \prod_{j=2}^n\frac{\sum_{a=1}^j e^{(2a-1)t}-1}{j-1}.
\end{equation}
For $n\ge2$, an exact factorization is
\begin{equation}\label{tfp:eq:minusfactor}
 \frac{P_n(-e^t)}{(n+1)!}
       =\frac{(-1)^n t}{n(n+1)}B_n(t).
\end{equation}
The factor $t$ comes from $1+q$. It is essential, not an error term.
Write
\begin{equation}\label{tfp:eq:lambda}
 \log B_n(t)=\nu_n t+\frac{U_nt^2}{2}
                  +\sum_{r\ge3}\lambda_{r,n}\frac{t^r}{r!}.
\end{equation}
The function $B_n$ is an analytic normalization, not a probability moment
generating function. In particular some individual signed variances are
negative. Nevertheless its aggregate quadratic coefficient is positive
for all sufficiently large $n$, and direct calculation gives
\begin{align}
 \nu_n&=\frac{n(n+3)}2-\frac32+H_{n-1},\label{tfp:eq:nu}\\
 U_n&=\frac{n^3}{9}-\frac{n^2}{6}-\frac{41n}{18}
               +\frac{29}{12}-3H_{n-1}-\HH{n-1}{2},\label{tfp:eq:U}\\
 \Delta_n:=\nu_n-\mu_n
       &=2n-\frac12-\frac1n-\frac1{n+1}.\label{tfp:eq:delta}
\end{align}
The large displacement $\Delta_n\sim2n$, though small relative to
$\sqrt{V_n}$, converts the odd Fourier factor into a nonzero local
contribution near the maximum.

For reproducibility, the signed moments for $j\ge2$ are
$\widetilde m_0(j)=1$ and
$\widetilde m_r(j)=\sum_{a=1}^j(2a-1)^r/(j-1)$ for $r\ge1$.
Use \eqref{tfp:eq:cumrec}, then add the cumulants of
$(e^t-1)/t$, the moment generating function of the continuous uniform
law on $[0,1]$. This generates all the $\lambda_{r,n}$.

\subsection{Fourier bounds with no hidden aperiodicity hypothesis}
Put $\Phi_n(\theta)=P_n(e^{\ii\theta})/(n+1)!$.
Fourier inversion gives
\begin{equation}\label{tfp:eq:fourier}
 \PP(S_n=k)=\frac1{2\pi}\int_{-\pi}^{\pi}
                    \Phi_n(\theta)e^{-\ii k\theta}\,\dd\theta.
\end{equation}
The two ends near $\pm\pi$ are one neighborhood on the circle.

\begin{lemma}[Two-neighborhood localization]\label{tfp:lem:arcs}
There are constants $a,c,C>0$ such that, for all large $n$,
\begin{align}
 |\Phi_n(\theta)|&\le e^{-cn^3\theta^2}
                      &&(|\theta|\le a/n),\label{tfp:eq:plusbound}\\
 |B_n(\ii\theta)|&\le C e^{-cn^3\theta^2}
                      &&(|\theta|\le a/n),\label{tfp:eq:minusbound}\\
 |\Phi_n(\theta)|&\le C e^{-cn}
                      &&(\dist(\theta,\pi\Z)\ge a/n).
                      \label{tfp:eq:minor}
\end{align}
The same bounds, with different constants, permit any fixed number of
Taylor or Fourier-differentiation operations used below.
\end{lemma}
\begin{proof}
A single characteristic factor satisfies
\begin{equation}\label{tfp:eq:factortrig}
 \frac{F_j(\ii\theta)}{j+1}
 =\frac{1+e^{\ii j\theta}\sin(j\theta)/\sin\theta}{j+1}.
\end{equation}
For $|\theta|\le a/n$, the identity
$|\E e^{\ii\theta X_j}|^2=\E\cos(\theta(X_j-X_j'))$, with $X_j'$ an
independent copy, and the inequality $1-\cos u\ge c_0u^2$ for small $u$
show that the product is at most
$\exp(-c_1\theta^2\sum_{j\ge j_0}j^2)$. This proves
\eqref{tfp:eq:plusbound}.

For the second neighborhood, the signed single-factor variance is
\begin{equation}
 \widetilde b_2(j)=\frac{j^2}{3}-\frac{2j}{3}-2
                     -\frac3{j-1}-\frac1{(j-1)^2}.
\end{equation}
It is at least a positive constant times $j^2$ for all sufficiently large
$j$. Uniform analyticity on $|t|\le a/j$ gives
\[
 \Re\log\!\left(
 \frac{\sum_{a=1}^j e^{(2a-1)\ii\theta}-1}{j-1}\right)
   =-\frac12\widetilde b_2(j)\theta^2+O(j^4\theta^4).
\]
Choose $a$ sufficiently small. The error is then absorbed into the
negative quadratic term. Finitely many small factors and
$(e^{\ii\theta}-1)/(\ii\theta)$ contribute only bounded factors,
proving \eqref{tfp:eq:minusbound}.

For completeness, set $\rho=\dist(\theta,\pi\Z)$. If
$\rho\ge K/n$ for a sufficiently large fixed $K$, then, for $j\ge n/2$,
\[
 \left|\frac{F_j(\ii\theta)}{j+1}\right|
 \le \frac{1+1/|\sin\theta|}{j+1}
 \le \frac2n+\frac{\pi}{K}<q<1.
\]
If instead $a/n\le\rho\le K/n$, put $u=n\rho$ and $v=j/n$.
On the compact rectangle $u\in[a,K]$, $v\in[1/2,1]$,
$|\sin(uv)/(uv)|$ has supremum strictly less than one. Formula
\eqref{tfp:eq:factortrig}, with its $O(1/j)$ extra term, gives the same
strict bound for all sufficiently large $n$. At least a positive
proportion of the factors thus have modulus at most $q<1$. The remaining
factors are characteristic functions and have modulus at most one.
This proves \eqref{tfp:eq:minor}. Fixed derivatives merely add polynomial
factors, which the exponential bounds absorb away from the centers;
inside the centers they are handled by the displayed analytic estimates.
\end{proof}

\begin{remark}[Exact parity balance does not remove the second peak]
Since $P_n(-1)=0$, the sums of the even- and odd-degree coefficients
are exactly equal. But \eqref{tfp:eq:minusfactor} has a simple zero, not an
identically vanishing neighborhood. A local coefficient expansion sees
its derivatives. Removing the first Bernoulli factor destroys the exact
parity balance and leaves a Fourier value of magnitude $2/[n(n+1)]$.
This is a concrete instance of the stability issue discussed in
\cite{dolgopyat}.
\end{remark}

\section{An arbitrary-order two-neighborhood expansion}\label{tfp:sec:allorders}
\subsection{Finite coefficient operators}
Let
\begin{equation}
 G_V(x)=\frac1{\sqrt{2\pi V}}e^{-x^2/(2V)},
 \qquad
 \He_r(x)=(-1)^r e^{x^2/2}\frac{\dd^r}{\dd x^r}e^{-x^2/2}.
\end{equation}
For a nonnegative integer $R$, use multi-indices
$\mathbf m=(m_3,\ldots,m_{R+2})$ with
\begin{equation}
 w(\mathbf m)=\sum_{r=3}^{R+2}(r-2)m_r\le R,
 \qquad D(\mathbf m)=\sum_{r=3}^{R+2}r m_r.
\end{equation}
An empty index set when $R=0$ means the single zero multi-index.
For a cumulant list $\boldsymbol\eta$, define
\begin{align}
 T_R(x;\boldsymbol\eta,V)
 &=\sum_{w(\mathbf m)\le R}
   \left[\prod_{r=3}^{R+2}
       \frac{\eta_r^{m_r}}{m_r!(r!)^{m_r}}\right]
       \frac{\He_{D(\mathbf m)}(x/\sqrt V)}{V^{D(\mathbf m)/2}},
       \label{tfp:eq:TR}\\
 S_R(x;\boldsymbol\eta,V)
 &=\sum_{w(\mathbf m)\le R}
   \left[\prod_{r=3}^{R+2}
       \frac{\eta_r^{m_r}}{m_r!(r!)^{m_r}}\right]
       \frac{\He_{D(\mathbf m)+1}(x/\sqrt V)}{V^{(D(\mathbf m)+1)/2}}.
       \label{tfp:eq:SR}
\end{align}
These are finite, directly computable expressions.

\begin{theorem}[All algebraic orders]\label{tfp:thm:allorders}
For every fixed nonnegative integer $R$, uniformly for all integer $k$,
\begin{equation}\label{tfp:eq:allorders}
\begin{split}
 \PP(S_n=k)={}&G_{V_n}(k-\mu_n)
          T_R(k-\mu_n;\boldsymbol\kappa_n,V_n)\\
 &+\frac{(-1)^{n+k}}{n(n+1)}G_{U_n}(k-\nu_n)
          S_R(k-\nu_n;\boldsymbol\lambda_n,U_n)\\
 &+O_R\!\left(V_n^{-1/2}n^{-(R+1)/2}\right).
\end{split}
\end{equation}
The assertion is an absolute-error estimate, not a relative estimate
at arbitrary tail indices. Once the two-candidate mode localization is
established, it gives the maximum to arbitrary relative algebraic order
by increasing $R$ and taking the maximum of the two candidate values.
\end{theorem}
\begin{proof}
Split \eqref{tfp:eq:fourier} into the neighborhoods from
Lemma \ref{tfp:lem:arcs}. The complement is exponentially small. At zero,
center the characteristic function and put $u=\theta\sqrt{V_n}$.
Its logarithm is
\begin{equation}\label{tfp:eq:localexp}
 -\frac{u^2}{2}
       +\sum_{r=3}^{R+2}\frac{\kappa_{r,n}}{r!V_n^{r/2}}(\ii u)^r
       +\text{remainder}.
\end{equation}
The coefficient for index $r$ is $O_r(n^{-(r-2)/2})$. Expanding the
exponential by this weight gives exactly the multi-indices in
\eqref{tfp:eq:TR}.

Here are the remainder details. Fix $R$, and first restrict
$|u|\le\sqrt{K\log n}$, where $K$ is a sufficiently large fixed number.
Uniform analytic cumulant bounds on $|t|\le a/n$ bound the logarithmic
Taylor remainder by
$C_R n^{-(R+1)/2}(1+|u|^{R+3})$ on that interval for all large $n$.
The omitted terms of the weighted exponential expansion have the same
order times a fixed polynomial in $|u|$, after multiplication by
$e^{-c u^2}$. Integration removes that polynomial without introducing
a power of $\log n$. Outside the cutoff, the Gaussian bounds make the
integral smaller than the required power by taking $K$ larger. This
proves an integrated error $O_R(n^{-(R+1)/2})$, before the Jacobian
$V_n^{-1/2}$.

Each retained integral is evaluated using
\begin{equation}\label{tfp:eq:hermiteintegral}
 \frac1{2\pi}\int_{\R}(\ii\theta)^r
       e^{-V\theta^2/2-\ii x\theta}\,\dd\theta
       =G_V(x)\frac{\He_r(x/\sqrt V)}{V^{r/2}}.
\end{equation}
Extending the integrals to the real line changes them only by the
already controlled Gaussian tails. All error estimates are uniform in
$x$ because $|e^{-\ii x\theta}|=1$.

At the other neighborhood, use \eqref{tfp:eq:minusfactor}. Fourier inversion
contributes $(-1)^k$, and the exact factor supplies $(-1)^n/[n(n+1)]$
and one additional $\ii\theta$. The same argument with $U_n$ and
$\lambda_{r,n}$ therefore gives \eqref{tfp:eq:SR}. Its own remainder is
$O_R(n^{-2}U_n^{-1}n^{-(R+1)/2})$, smaller than the displayed total
remainder. This proves \eqref{tfp:eq:allorders}.
\end{proof}

\subsection{What ``all orders'' means here}
The theorem is not a claim of convergence of an infinite asymptotic
series. For every desired algebraic error, it supplies a finite formula
and a proved remainder. Using \eqref{tfp:eq:cumrec}, polynomial division,
Faulhaber's formulas, and Euler--Maclaurin expansion of harmonic numbers
turns the finite formula into powers of $n^{-1}$ with logarithmic
coefficients and bounded lattice variables. Constants can involve
$\gamma$ and values $\zeta(r)$.

For bounded $k-\mu_n$, half-integer powers in the standardized Hermite
notation combine into integer powers: odd Hermite polynomials supply
an additional odd factor of their small argument. The second
neighborhood has the same property, with $k-\nu_n=-2n+O(1)$.
Taking a maximum preserves an error bound but need not preserve a
single smooth coefficient function across a mode transition.
We establish neither exponentially improved sectors nor Borel
summability. Those stronger meanings of ``full transseries'' are left
open.

\section{Derivation of the explicit coefficients}\label{tfp:sec:coeff}
\subsection{The contribution from zero}
Use Lemma \ref{tfp:lem:cumulants} in \eqref{tfp:eq:allorders}. Up to relative
order $n^{-4}$ at bounded $d=k-\mu_n$, the only odd-cumulant term needed
is $-\kappa_{3,n}d/(2V_n^2)$. Even terms are generated by the partitions
of weights at most four using $\kappa_4,\kappa_6,\kappa_8,\kappa_{10}$,
where $\kappa_{2r}$ has weight $r-1$. As examples, the central factors
through weight three are
\begin{equation}\label{tfp:eq:centralfactors}
\begin{split}
 1+\frac{\kappa_4}{8V^2}
 -\frac{\kappa_6}{48V^3}
 +\frac{35\kappa_4^2}{384V^4}
 +\frac{\kappa_8}{384V^4}
 -\frac{7\kappa_4\kappa_6}{128V^5}
 +\frac{385\kappa_4^3}{3072V^6}.
\end{split}
\end{equation}
Here and in the following calculation the subscript $n$ is suppressed.
For nonzero $d$, use the full Hermite polynomial, not merely its
constant term.

The relevant polynomial portions of the cumulants begin as follows:
\begin{align}
 \kappa_4={}&-\frac{2n^5}{75}-\frac{n^4}{5}
        -\frac{26n^3}{45}+\frac{26n^2}{5}+O(n),\label{tfp:eq:k4poly}\\
 \kappa_6={}&\frac{16n^7}{441}+\frac{8n^6}{21}
        +\frac{104n^5}{63}+O(n^4),\label{tfp:eq:k6poly}\\
 \kappa_8={}&-\frac{16n^9}{135}-\frac{8n^8}{5}+O(n^7),\\
 \kappa_{10}={}&\frac{256n^{11}}{363}+O(n^{10}).
\end{align}
These formulas follow from \eqref{tfp:eq:cumrec}; the accompanying symbolic
script performs the exact rational operations.

The factorial and variance factors are
\begin{align}
 \frac{(n+1)!}{\sqrt{2\pi}\,n^{n+3/2}e^{-n}}
  &=(1+n^{-1})\exp\!\left(\frac1{12n}-\frac1{360n^3}
                           +O(n^{-5})\right),\label{tfp:eq:stirling}\\
 \frac{n^{3/2}}{3\sqrt V}
  &=\left(1+\frac9{2n}-\frac{29}{2n^2}
                +\frac{27h_n-9h_{2,n}-18}{n^3}\right)^{-1/2}.
                \label{tfp:eq:varfactor}
\end{align}
Multiplying the Hermite contributions, $e^{-d^2/(2V)}$, and these two
factors gives $1+a_1/n+a_2/n^2+C_3/n^3+C_4/n^4$.
For a quick check on the first coefficient,
\begin{equation}
 a_1=\underbrace{\frac{13}{12}}_{\text{factorial}}
          -\underbrace{\frac94}_{\text{variance}}
          -\underbrace{\frac{27}{100}}_{\text{fourth cumulant}}
        =-\frac{431}{300}.
\end{equation}
The first lattice loss is $-d^2/(2V)\sim-9d^2/(2n^3)$,
and the first skew displacement term is
$-\kappa_3d/(2V^2)\sim-81d/(2n^4)$. These are useful independent checks
on \eqref{tfp:eq:C3} and \eqref{tfp:eq:C4}.

\subsection{The contribution from \texorpdfstring{$\pi$}{pi}}
For bounded $d$ and $x=k-\nu_n=d-\Delta_n$, the leading term from the
second neighborhood is
\begin{equation}\label{tfp:eq:piapprox}
 \frac{(-1)^{n+k}}{n(n+1)}G_{U_n}(x)\frac{x}{U_n}.
\end{equation}
Now $x=-2n+O(1)$, $U_n\sim n^3/9$, and
$G_{U_n}(x)/G_{V_n}(0)=1+O(n^{-1})$. Therefore
\begin{equation}\label{tfp:eq:pileading}
 \frac{\text{second-neighborhood contribution}}{G_{V_n}(0)}
   =-\frac{18(-1)^{n+k}}{n^4}+O(n^{-5}).
\end{equation}
The quartic correction changes this only at the next order, and the
odd-cumulant bounds in Lemma \ref{tfp:lem:cumulants} control the remaining
terms. The normalization from \eqref{tfp:eq:stirling} and
\eqref{tfp:eq:varfactor} is $1+O(n^{-1})$, so the same coefficient $-18$
appears in \eqref{tfp:eq:fourorders}. Taylor expansion of the remaining
explicit factors, keeping $h_n$ and $h_{2,n}$ exact, gives a uniform
remainder $O_D((1+\log^2 n)n^{-5})$. This proves the local expansion
part of Theorem \ref{tfp:thm:main}.

\section{The maximizing exponent and infinitely many rounding failures}
\label{tfp:sec:modes}
\subsection{Global localization of the mode}
A local expansion at a proposed maximum does not by itself prove that
this is the global maximum. We supply that step separately.

\begin{lemma}[Principal smooth peak and secondary bound]\label{tfp:lem:smooth}
Let $A_n(d)$ be the real-valued contribution of $|\theta|\le a/n$ to
\eqref{tfp:eq:fourier}, centered so that $d=k-\mu_n$. It is defined for real
$d$. For a sufficiently small fixed $\varepsilon>0$, it has a unique
critical point $d_n^*$ in $|d|\le\varepsilon\sqrt{V_n}$, and
\begin{equation}\label{tfp:eq:smoothmode}
 d_n^*=-\frac{\kappa_{3,n}}{2V_n}+O(n^{-2})
       =-\frac9{2n}+O(n^{-2}).
\end{equation}
On that interval $A_n$ is strictly concave, with curvature comparable
to $-G_{V_n}(0)/V_n$. If $B_n^{\rm arc}(d)$ denotes the other
neighborhood's contribution, including its parity factor at integer
indices, then
\begin{equation}\label{tfp:eq:secondarybound}
 |B_n^{\rm arc}(d)|\le
 C G_{V_n}(0)\frac{|d|+n}{n^5}
 \qquad (|d|\le\varepsilon\sqrt{V_n}).
\end{equation}
\end{lemma}
\begin{proof}
Differentiate the principal Fourier integral. The proof of
Theorem \ref{tfp:thm:allorders} applies with additional powers of
$\theta$. Uniformly on the stated interval, its second derivative is
\[
 A_n''(d)=-\frac{G_{V_n}(0)}{V_n}
               (1+O(\varepsilon^2+n^{-1})).
\]
At zero, even cumulants contribute nothing to the first derivative.
The cubic term and Lemma \ref{tfp:lem:cumulants} give
\[
 A_n'(0)=G_{V_n}(0)
       \left(-\frac{\kappa_{3,n}}{2V_n^2}+O(n^{-5})\right).
\]
Strict concavity and the intermediate value theorem produce the unique
critical point and \eqref{tfp:eq:smoothmode}.

For the other arc, the extra $\ii\theta$ differentiates its smooth
Gaussian profile. Its even-cumulant terms are odd in $x=d-\Delta_n$
and bounded, on a fixed standardized interval, by
$C G_{V_n}(0)|x|/U_n$. The odd-cumulant terms contribute at most
$C G_{V_n}(0)n^{-4}$ by Lemma \ref{tfp:lem:cumulants}. Multiplication by
$1/[n(n+1)]$, use of $|x|\le|d|+Cn$, and sufficiently many terms of
\eqref{tfp:eq:allorders} give \eqref{tfp:eq:secondarybound}.
\end{proof}

\begin{proposition}[Two integer candidates]\label{tfp:prop:twocandidates}
For all sufficiently large $n$, every global maximizing exponent lies
in $\{\lfloor\mu_n\rfloor,\lceil\mu_n\rceil\}$.
\end{proposition}
\begin{proof}
The uniform absolute-error expansion in Theorem \ref{tfp:thm:allorders}
first implies a uniform local central limit estimate with error
$o(G_{V_n}(0))$. Comparing with an integer at distance at most $1/2$
from $\mu_n$ confines every maximizer to
$|d|\le\varepsilon\sqrt{V_n}$, for any fixed small $\varepsilon$ and
all large $n$.

Within this interval, compare to the integer nearest $\mu_n$.
Any integer outside the two stated candidates has $|d|\ge1$, while
the comparison integer has distance at most $1/2$. Lemma
\ref{tfp:lem:smooth}, with $\varepsilon$ chosen small enough that the
upper and lower curvature bounds have ratio arbitrarily close to one,
shows a principal-arc loss at least
\[
 cG_{V_n}(0)\frac{1+d^2}{n^3}
 \quad\text{for such a noncandidate.}
\]
The displacement $d_n^*=O(n^{-1})$ does not alter this fixed quadratic
gap. The total secondary-arc advantage is at most
$CG_{V_n}(0)(|d|+n)n^{-5}$, which is uniformly smaller. The minor-arc
error is exponentially small. Thus a noncandidate cannot maximize.
\end{proof}

Proposition \ref{tfp:prop:twocandidates} completes the proof of
Theorem \ref{tfp:thm:main}: both candidate offsets are bounded, and taking
a maximum is Lipschitz in the sup norm of the two errors.

\subsection{The parity-dependent transition}
Put
\begin{equation}
 m=\lfloor\mu_n\rfloor,\qquad
 r_n=\mu_n-m-\frac12,\qquad s_n=(-1)^{n+m}.
\end{equation}
Subtract \eqref{tfp:eq:fourorders} at $m$ and $m+1$. The harmonic and
constant terms cancel. Since $2(m-\mu_n)+1=-2r_n$,
\begin{equation}\label{tfp:eq:diffexact}
 \frac{c_{n,m+1}-c_{n,m}}{\calE_n}
  =\frac{9r_n}{n^3}
   +\frac{-\frac{1263}{20}r_n-\frac{81}{2}+36s_n}{n^4}
   +O\!\left(\frac{1+\log^2 n}{n^5}\right).
\end{equation}
In the transition window $r_n=O(n^{-1})$, this becomes
\begin{equation}\label{tfp:eq:threshold}
 \frac{c_{n,m+1}-c_{n,m}}{\calE_n}
 =\frac9{n^3}\left(r_n-\frac{\frac92-4s_n}{n}
                +O\!\left(\frac{1+\log^2n}{n^2}\right)\right).
\end{equation}
Thus the transition is at $r_n\approx1/(2n)$ when $s_n=1$, but at
$r_n\approx17/(2n)$ when $s_n=-1$. Rounding the principal smooth mode
alone would instead put it at $9/(2n)$ for both parities. Formula
\eqref{tfp:eq:threshold} decides the sign whenever the separation from
the proposed boundary dominates its displayed error. It is not an
exact finite-$n$ rounding rule inside the unresolved boundary layer.

\begin{theorem}[The smaller Fourier contribution changes the mode
infinitely often]\label{tfp:thm:failures}
Let $\widehat k_n$ be a nearest integer to
$\mu_n-9/(2n)$, with any fixed tie convention. Then $\widehat k_n$ fails
to maximize $c_{n,k}$ for infinitely many $n$. The same assertion holds
for rounding the exact principal smooth peak
$\mu_n+d_n^*$.
If $E(N)$ counts failures of the explicit rule $\widehat k_n$ up to $N$,
then, for all sufficiently large $N$,
\begin{equation}\label{tfp:eq:failurecount}
 \log N-O(1)\le E(N)\le12\log N+O(1).
\end{equation}
In particular the failures have order $\log N$ and natural density zero.
\end{theorem}
\begin{proof}
Let $b_n=H_{n+1}-1$ and, for each sufficiently large integer $j$, let
$n_j$ be the first index with $b_{n_j}>j+1/2$.
The overshoot is between zero and $1/(n_j+1)$ because
$b_n-b_{n-1}=1/(n+1)$. For $n=n_j+\ell$, $\ell=1,2,3$, therefore
\begin{equation}\label{tfp:eq:overshoot}
 \frac{1-o(1)}n\le r_n\le\frac{4+o(1)}n,
 \qquad \lfloor b_n\rfloor=j.
\end{equation}
Also
\[
 s_n=(-1)^{n+\lfloor\mu_n\rfloor}
     =(-1)^{n(n+1)/2+j}.
\]
The parity of $n(n+1)/2$ has period four, with two occurrences of each
parity in each period. In any three consecutive indices there is an
$s_n=1$. At such an index in \eqref{tfp:eq:overshoot}, the true transition
boundary $1/(2n)$ lies below $r_n$ by a fixed multiple of $1/n$,
whereas the principal boundary $9/(2n)$ lies above it by a fixed
multiple of $1/n$. The true maximizing index is $m+1$ and the
principal rounding rule chooses $m$. All margins exceed the error in
\eqref{tfp:eq:threshold}; the $O(n^{-2})$ correction in
\eqref{tfp:eq:smoothmode} also cannot change this conclusion.

There are $\log N+O(1)$ crossings up to $N$, since
$b_N=\log N+\gamma-1+o(1)$, and the three-index windows are eventually
disjoint. This proves the lower bound, allowing an $O(1)$ loss for
the last window.

For the upper bound, \eqref{tfp:eq:diffexact} and
\eqref{tfp:eq:threshold} imply that a failure is possible only when
$0<r_n<10/n$, for all sufficiently large $n$; outside this interval
both rules choose the same candidate. Such an index lies just after a
half-integer crossing. If it is $n=n_j+\ell$, then
$r_n\ge\ell/(n+1)$, so $r_n<10/n$ permits at most twelve indices per
crossing. Finitely many earlier indices are absorbed into $O(1)$.
\end{proof}

\begin{remark}
Rounding the mean itself also fails infinitely often. In the first
three indices at or immediately after a half-integer crossing, select
$s_n=-1$. The overshoot is at most $(3+o(1))/n$, below the true
$17/(2n)$ boundary, even though the nearest integer to the mean is
already the upper one. No equidistribution assumption about harmonic
numbers or exponential fractional parts is used in either argument.
\end{remark}

\section{The log-periodic term and a limitation of power--log expansions}
\label{tfp:sec:logperiodic}
Set $\delta(x)=\dist(x,\Z)$ and
\begin{equation}
 K_0=-\frac{27}{2}\gamma+\frac92\zeta(2)
             -\frac{107751433393}{1587600000}.
\end{equation}
The nearest-integer loss at third order is independent of which of
two almost tied candidates eventually wins. Indeed, a departure from
the nearest-mean candidate occurs only within $O(1/n)$ of a tie and
changes the squared distance by $O(1/n)$.
Using $h_n=\log n+\gamma+O(1/n)$ and
$h_{2,n}=\zeta(2)+O(1/n)$ in Theorem \ref{tfp:thm:main} gives
\begin{corollary}[Explicit log-periodic correction]\label{tfp:cor:periodic}
\begin{equation}\label{tfp:eq:periodic}
\begin{split}
 \frac{M_n}{\calE_n}
 =1+\frac{a_1}{n}+\frac{a_2}{n^2}
 +\frac{-\frac{27}{2}\log n+K_0
          -\frac92\delta(\log n+\gamma-1)^2}{n^3}
 +O\!\left(\frac{\log n}{n^4}\right).
\end{split}
\end{equation}
There is no polynomial $Q$ for which the same expression can be
replaced by
$1+a_1/n+a_2/n^2+Q(\log n)/n^3+o(n^{-3})$.
\end{corollary}
\begin{proof}
The distance-to-integers function, and its square on the bounded range
$[0,1/2]$, are Lipschitz. Replacing the exact harmonic argument by
$\log n+\gamma-1$ therefore changes the third-order term by $O(n^{-4})$.
This proves \eqref{tfp:eq:periodic}.

For the second assertion, a possible polynomial would have to equal
$-(27/2)x$ plus a constant and an $o(1)$ error after the bounded
remainder is isolated. But for $a=0$ and $a=1/2$, integers $n_j$ nearest
$\exp(j+1-\gamma+a)$ satisfy
$\log n_j+\gamma-1-j\to a$. The squared distance has the two distinct
limits $0$ and $1/4$. A single constant cannot absorb both.
\end{proof}

The periodic term has the absolutely convergent Fourier representation
\begin{equation}\label{tfp:eq:periodicfourier}
 \delta(x)^2=\frac1{12}
       +\frac1{\pi^2}\sum_{\ell=1}^{\infty}
                  \frac{(-1)^\ell}{\ell^2}\cos(2\pi\ell x).
\end{equation}
It follows by integrating twice to compute the Fourier coefficients of
$x^2$ on $[-1/2,1/2]$. Thus the third-order correction can alternatively
be represented using complex powers $n^{2\pi\ii\ell}$. It is periodic
in $\log n$ with period one, or multiplicatively periodic at leading
order under $n\mapsto\ee n$. This oscillation is generated by taking a
lattice maximum, not by a singularity of the original probability
model. The additional parity term at fourth order has a different
origin, namely the second Fourier neighborhood.

\section{A fixed-smoothing hierarchy}\label{tfp:sec:smoothing}
The role of the simple zero at $-1$ can be varied without changing the
growing part of the model. For a fixed integer $r\ge0$, define
\begin{equation}\label{tfp:eq:Pr}
 P_{n,r}(q)=(1+q)^r\prod_{j=2}^n
              \left(1+\sum_{a=1}^j q^{2a-1}\right).
\end{equation}
Its total mass is $2^{r-1}(n+1)!$. Let $\mu_{n,r}$ and $V_{n,r}$ be the
mean and variance after normalization. Then
\begin{equation}
 \mu_{n,r}=\mu_n+\frac{r-1}{2},\qquad
 V_{n,r}=V_n+\frac{r-1}{4}.
\end{equation}
The original sequence is $r=1$.

\begin{theorem}[Zero multiplicity and the second-neighborhood order]
\label{tfp:thm:smoothing}
For fixed $r\ge0$ and bounded $k-\mu_{n,r}$, the contribution of the
$\pi$ neighborhood to the normalized coefficient probability, divided
by $G_{V_{n,r}}(0)$, is
\begin{equation}\label{tfp:eq:smoothinglaw}
 (-1)^{n-1+r+k}
 \begin{cases}
  \displaystyle 2^{1-r}3^r(-1)^{r/2}(r-1)!!\,
                  n^{-2-3r/2}(1+O(n^{-1})),&r\text{ even},\\[5pt]
  \displaystyle 2^{2-r}3^{r+1}(-1)^{(r+1)/2}r!!\,
                  n^{-(3r+5)/2}(1+O(n^{-1})),&r\text{ odd}.
 \end{cases}
\end{equation}
The convention is $(-1)!!=1$. In particular the first orders are
$n^{-2}$ for $r=0$, $n^{-4}$ for $r=1$, and $n^{-5}$ for $r=2$.
\end{theorem}
\begin{proof}
The exact analogue of \eqref{tfp:eq:minusfactor} is
\begin{equation}
 \frac{P_{n,r}(-e^t)}{2^{r-1}(n+1)!}
 =\frac{(-1)^{n-1+r}2^{1-r}t^r}{n(n+1)}B_{n,r}(t),
\end{equation}
where
\[
 B_{n,r}(t)=\left(\frac{e^t-1}{t}\right)^r
          \prod_{j=2}^n\frac{\sum_{a=1}^j e^{(2a-1)t}-1}{j-1}.
\]
Its center is $\nu_n+(r-1)/2$, its quadratic coefficient is
$U_n+(r-1)/12$, and the displacement from the principal center is still
$\Delta_n$. The $t^r$ factor gives the $r$th Hermite derivative of the
Gaussian. At $x=-2n+O(1)$ the standardized argument tends to zero at
rate $n^{-1/2}$.

For even $r$, use
$\He_r(0)=(-1)^{r/2}(r-1)!!$ and
$U^{-r/2}\sim3^rn^{-3r/2}$. For odd $r$, use
$\He_r(z)=(-1)^{(r-1)/2}r!!\,z+O(z^3)$, with
$z=x/\sqrt U\sim-6n^{-1/2}$. Combine with
$2^{1-r}/[n(n+1)]$. The Gaussian ratio is $1+O(n^{-1})$.
Even-cumulant corrections have relative order $n^{-1}$ or smaller;
the odd-cumulant cancellation from Lemma \ref{tfp:lem:cumulants} makes
those corrections smaller still. This proves \eqref{tfp:eq:smoothinglaw}.
\end{proof}

For $r=0$ the secondary contribution is larger than the ordinary
$n^{-3}$ rounding loss. Accordingly the two-candidate mode theorem for
$r=1$ must not be transferred to $r=0$ without a new argument. For
$r\ge2$ the parity effect occurs later. The theorem isolates a general
principle: a finite smoothing factor can move a Fourier contribution
across several orders, and odd zero multiplicity has an additional
cancellation near the moving center.

\section{Inverse growth and the integer threshold}\label{tfp:sec:inverse}
\subsection{A logarithmic core}
From \eqref{tfp:eq:leading}, taking logarithms gives
\begin{equation}\label{tfp:eq:logM}
 \log(M_n/3)
 =n(\log n-1)+\frac{A}{n}+\frac{B}{n^2}
                    +O\!\left(\frac{\log n}{n^3}\right),
 \quad A=-\frac{431}{300},\quad B=\frac{5907087}{490000}.
\end{equation}
The second constant is $B=a_2-a_1^2/2$; confusing a multiplicative
coefficient with a logarithmic coefficient would give the wrong inverse.

For large $y$, let
\begin{equation}\label{tfp:eq:lambert}
 L=\log(y/3),\qquad x=\frac{L}{W_0(L/\ee)}.
\end{equation}
The principal real Lambert function is characterized by
$W_0(z)e^{W_0(z)}=z$. It follows directly that
$x(\log x-1)=L$. This is the same core-and-correction organization used
in the repository's inversion materials \cite{transseries}.
[Added October 1, 2026, batch 73O2: in the canonical volume
\emph{Transseries and Inversion}~\cite{tai} the exact solution of this core
is its theorem ``Exact solution of the dominant block''
(\texttt{p0:thm:lambert-core}), and the two corrections of
Theorem~\ref{tfp:thm:inverse} are an instance of its theorem ``Perturbed
inversion'' (\texttt{p0:thm:perturbed-inversion}). The manuscript re-derives
them for its own forward map, which is a pure power carrier $(n/\ee)^n$
rather than a gamma function. No novelty is claimed for the inversion
method; what is specific to this report is the forward input $A$, $B$ and the
lattice-dependent third coefficient \eqref{tfp:eq:logthird}.]

\begin{theorem}[Inverse expansion on the sequence values]\label{tfp:thm:inverse}
When $y=M_n$ and $x$ is given by \eqref{tfp:eq:lambert},
\begin{equation}\label{tfp:eq:inverse}
 n=x+\frac{431}{300x\log x}
       -\frac{5907087}{490000x^2\log x}+O(x^{-3}).
\end{equation}
The error concerns the recovered real index at the exact sequence
values. It is not an $o(1)$ approximation to a staircase at arbitrary
arguments without rounding.
\end{theorem}
\begin{proof}
Let $F(t)=t(\log t-1)$, so $F'(t)=\log t$. Equation \eqref{tfp:eq:logM}
first gives $n-x=O(1/(x\log x))$ by the mean value theorem.
Write $n=x+\varepsilon$ and expand:
\[
 0=(\log x)\varepsilon+\frac A x+\frac B{x^2}
        +O\!\left(\frac{\log x}{x^3}\right).
\]
Here the quadratic Taylor term and the changes in $A/n$ and $B/n^2$
are smaller than the stated remainder. Division by $\log x$ gives
\eqref{tfp:eq:inverse}.
\end{proof}

The next logarithmic coefficient, retaining a chosen candidate offset
$d$ and exact harmonic quantities, is
\begin{equation}\label{tfp:eq:logthird}
 -\frac92d^2-\frac{27}2h_n+\frac92h_{2,n}
                        -\frac{197099473}{3937500}.
\end{equation}
Thus further inversion is possible, but the lattice dependence must be
retained or treated by candidate envelopes. It cannot be replaced by
an invented smooth third coefficient.

\subsection{Thresholds and a rigorous rounding qualification}
The coefficients are nonnegative, and the next factor has constant
coefficient one, so $M_{n+1}\ge M_n$. Formula \eqref{tfp:eq:leading} also
gives $M_{n+1}/M_n\sim n$, hence strict increase eventually. Define
\begin{equation}
 N(y)=\min\{n\ge1:M_n\ge y\}.
\end{equation}
Let $t(y)$ be the large real solution of
\begin{equation}\label{tfp:eq:smoothinverse}
 F(t)+A/t+B/t^2=\log(y/3).
\end{equation}
The left side is strictly increasing for large $t$.
There is a constant $C>0$ such that, for all sufficiently large $y$,
\begin{equation}\label{tfp:eq:thresholdbracket}
 \left\lceil t(y)-Ct(y)^{-3}\right\rceil
       \le N(y)\le
 \left\lceil t(y)+Ct(y)^{-3}\right\rceil.
\end{equation}
To prove this, transport the $O(\log n/n^3)$ error in \eqref{tfp:eq:logM}
through a derivative bounded below by a positive multiple of $\log n$.
An index smaller than $t-Ct^{-3}$ cannot reach $y$, while an integer
at least $t+Ct^{-3}$ must reach it. Consequently $N(y)=\lceil t(y)\rceil$
away from the indicated small neighborhoods of integers. No exact
ceiling rule is asserted in the boundary layer.
[Added October 1, 2026, batch 73O2: this bracket is an instance of the
volume's staircase theorem (\texttt{p0:thm:staircase}~\cite{tai}), whose
part~(1) gives $N(y)=\lceil\nu(y)\rceil$ for every admissible increasing
interpolation $\nu$ and whose part~(2), the separation condition, decides the
ceiling from an approximation $x$ with $|x-\nu(y)|\le\eta$ exactly when
$\lceil x-\eta\rceil=\lceil x+\eta\rceil$. The separation step and its
failure for arbitrarily small $\eta$ are formalized in Lean as
\texttt{Fabius.staircase\_separation} and
\texttt{Fabius.staircase\_separation\_fails} in
\path{Analysis/FabiusFunction/Lean/FabiusFunction/StaircaseInversion.lean}.
Those declarations are general; they do not verify any statement of this
report.]

For a smooth approximation $\widetilde M_J(t)$ with relative error
$O(t^{-J})$, the same argument gives an index error
$O(t^{-J}/\log t)$ on matched values, provided its logarithmic derivative
is comparable to $\log t$. For the all-order result here, this principle
is applied to the two lattice-resolved candidate expressions rather
than to a nonexistent globally smooth power--log series for the maximum.

\section{Exact computation and reproducibility}\label{tfp:sec:computation}
\subsection{A quadratic-work convolution step}
Let $a_k=c_{n-1,k}$, extended by zero. Multiplication by the $n$th factor
requires
\begin{equation}
 c_{n,k}=a_k+s_k,\qquad
 s_k=\sum_{j=1}^n a_{k-(2j-1)}.
\end{equation}
The sliding recurrence
\begin{equation}\label{tfp:eq:sliding}
 s_k=s_{k-2}+a_{k-1}-a_{k-2n-1}
\end{equation}
uses separate running sums for the two parities. It computes a whole
row in $O(n^2)$ integer additions and subtractions, hence rows through
$N$ in $O(N^3)$ such operations, with $O(N^2)$ live integer entries.
This is an arithmetic-operation count, not a bit-complexity bound:
the integers grow with $N$.

The supplied program computes every row through $n=400$ using arbitrary-
precision integers, checks total mass against the exact integer
$(n+1)!$, and matches the first twenty maximum values printed in
A039831 \cite{oeis831}. The record's b-file was not needed for these
calculations. An independent naive convolution on small rows and exact
moment identities are included as unit tests. Symbolic checks derive
the four coefficients independently from the cumulant recurrence and
Hermite integration. These finite checks support implementation
correctness; they do not replace the analytic proofs.

\subsection{Accuracy of the expansion}
The table reports $(\text{approximation}-M_n)/M_n$. ``Order $j$'' means
\eqref{tfp:eq:fourorders} through $n^{-j}$ at the exactly computed maximizing
exponent, including the second neighborhood for order four.
\begin{center}
\small
\begin{tabular}{r r r r r}
\toprule
$n$ & Leading only & Order two & Order three & Order four\\
\midrule
50  & $2.497\!\times10^{-2}$ & $8.835\!\times10^{-4}$ & $-1.218\!\times10^{-4}$ & $2.010\!\times10^{-5}$\\
100 & $1.336\!\times10^{-2}$ & $1.240\!\times10^{-4}$ & $-8.562\!\times10^{-6}$ & $7.063\!\times10^{-7}$\\
200 & $6.921\!\times10^{-3}$ & $1.701\!\times10^{-5}$ & $-6.047\!\times10^{-7}$ & $2.581\!\times10^{-8}$\\
400 & $3.525\!\times10^{-3}$ & $2.313\!\times10^{-6}$ & $-3.933\!\times10^{-8}$ & $8.016\!\times10^{-10}$\\
\bottomrule
\end{tabular}
\end{center}
The computations use high-precision floating point only after exact
integer coefficients have been obtained. The displayed errors are
numerical measurements, not outward-rounded interval bounds.

The separate program \texttt{all\_orders.py} evaluates the finite
formula of Theorem \ref{tfp:thm:allorders} directly, with exact rational
cumulants before numerical Gaussian--Hermite evaluation. At $n=100$,
$k=4954$, and weight budget $R=8$, its relative error is about
$1.146\times10^{-10}$. Retaining exact cumulants and variance in this
finite expression automatically retains some higher-order terms that
are absent from the four-power truncation. Increasing the order at a
fixed $n$ is not guaranteed to improve an asymptotic approximation.

Examples in the exact calculation where rounding
$\mu_n-9/(2n)$ chooses the wrong maximizing exponent include
$n=139,140,141,142$ and $n=373,374,379,380$. At $n=139$,
$\mu_n\approx9595.522425$: principal rounding selects $9595$, but the
actual unique maximizing exponent is $9596$.
The first-order parity threshold itself is asymptotic: for instance its
unqualified prediction at $n=56$ is wrong. This does not contradict
\eqref{tfp:eq:threshold}, which retains a boundary-layer error and asserts
no universal small-$n$ rule.

\subsection{Inverse-index errors}
At the exact input $y=M_n$, the errors in the Lambert core and in
\eqref{tfp:eq:inverse} are:
\begin{center}
\begin{tabular}{r r r}
\toprule
$n$ & $x-n$ & Two-correction inverse minus $n$\\
\midrule
50  & $-6.3043\times10^{-3}$ & $-1.9121\times10^{-4}$\\
100 & $-2.8811\times10^{-3}$ & $-2.3115\times10^{-5}$\\
200 & $-1.3017\times10^{-3}$ & $-2.7890\times10^{-6}$\\
400 & $-5.8723\times10^{-4}$ & $-3.3910\times10^{-7}$\\
\bottomrule
\end{tabular}
\end{center}
The complete integer maxima, modes, means, relative errors, and inverse
checks are supplied in CSV files. Reproduction requires no network
access once the package's Python dependencies are installed.

\section{Further research questions and topics}\label{tfp:sec:questions}
The following problems are not claimed solved by this article.

\paragraph{1. Effective thresholds and interval certificates.}
Make the constants in the two-neighborhood bounds explicit, prove a
numerical threshold for Proposition \ref{tfp:prop:twocandidates}, and join
it to exhaustive exact verification below that threshold. This would
turn an eventual theorem into a fully certified mode algorithm for all
$n$. The main difficulty is obtaining useful, not merely enormous,
constants in the analytic remainder.

\paragraph{2. Exact mode-transition counting.}
Theorem \ref{tfp:thm:failures} proves only upper and lower logarithmic
bounds. Is there a constant $c$ with $E(N)\sim c\log N$? The leading
four-periodic picture suggests four failures near a generic
half-integer harmonic crossing, but near-boundary overshoots require
control before this can be a theorem. Determining whether those
exceptional overshoots affect the leading counting constant is a
natural quantitative problem.

\paragraph{3. The distinct-coefficient conjecture A039824.}
Resolve the conjectured formula $n^2-3$ for all $n>6$
\cite{oeis824}. The present central analysis can constrain nearby equal
coefficients but does not rule out distant coincidences across the
whole row. A proof would need additional exact algebra or a global
coefficient-ordering argument. This problem must not be conflated with
asymptotic knowledge of the largest coefficient.

\paragraph{4. General fixed smoothing polynomials.}
Replace $(1+q)^r$ in \eqref{tfp:eq:Pr} by a fixed polynomial $Q(q)$ with
nonnegative coefficients. Determine a complete expansion in terms of
the order of vanishing and the initial Taylor coefficients of $Q$ at
roots of unity. Theorem \ref{tfp:thm:smoothing} gives the simplest
one-root instance, including its odd-multiplicity cancellation.

\paragraph{5. Growing smoothing.}
Let $r=r(n)$ grow. Determine the regimes in which suppression of the
secondary neighborhood competes with the change in the principal
variance and with the lattice spacing. The proof for fixed $r$ is not
uniform in $r$ and cannot be reused without new estimates.

\paragraph{6. More residue classes and several resonant neighborhoods.}
For products supported on zero plus a growing arithmetic progression,
classify which roots of unity contribute at algebraic orders. Develop
a mode-selection rule that compares several residue classes rather
than only two parities. A useful objective is a finite-state
lattice-envelope calculus for asymptotic maxima.

\paragraph{7. Exponentially improved expansions.}
After the two algebraic neighborhoods have been expanded to high order,
identify the next exponential scales, optimal truncation indices, and
possible Stokes effects under complex deformation. The exponentially
small minor-arc bound in this article is not an identification of those
scales and is not a resurgent description.

\paragraph{8. Formal verification of the analytic-to-discrete bridge.}
Formalize the product normalization, cumulant recurrence, Fourier
integral, and two-neighborhood estimates, then verify the coefficient
arithmetic and mode comparison. The most reusable component would be
a theorem transporting uniform local expansions into a finite
candidate set for a discrete maximum. Such a theorem would apply
beyond this particular OEIS entry.

\paragraph{9. Certified inverse thresholds with lattice data.}
Build an inverse algorithm that uses the Lambert core to find a small
candidate set of indices and evaluates rigorously bounded two-sector
formulas there. Prove exact decisions except in explicitly marked
intervals, falling back to exact coefficient computation when needed.
This would implement the distinction between a real inverse
approximation and an integer staircase in Section \ref{tfp:sec:inverse}.
[Added October 1, 2026, batch 73O2: the candidate-set step is the volume's
proposition ``Certified discrete inversion''
(\texttt{p0:prop:certified-rounding}~\cite{tai}); what remains open is the
rigorous two-sector bound that feeds it.]

\section{Proof status and suggested integration into ProveIt}
\subsection{What has been established}
The paper proves the conjectured leading asymptotic for A039831, the
uniform four-order local expansion, an arbitrary-order coefficient
operator with a remainder, eventual two-candidate localization,
parity-dependent transitions, infinitely many failures of principal
rounding with logarithmic counting bounds, the nonexistence of a pure
power--log third-order coefficient for the maximum, the fixed-smoothing
law, and the displayed inverse-growth and threshold statements.

These are ordinary mathematical proofs in the article. They have not
been checked by Lean, Rocq, or an independent human referee in this
work. The exact computations and symbolic assertions are executable
checks of finite claims, not a formalization of the asymptotic theorems.
The phrases ``proved here'' and ``current OEIS conjecture'' have these
separate meanings throughout.

\subsection{A reusable formalization order}
A suitable first layer would define the finite coefficient product and
prove the sliding recurrence, total mass, and probability
normalization. A second layer would formalize the raw-moment/cumulant
recurrence and the exact formulas for $\mu_n,V_n,\kappa_{3,n}$.
The symbolic identities in the supplied script are natural candidates
for exact rational reflection certificates.

The analytic layer would prove the two-neighborhood localization and
Hermite integral identity, followed by a fixed-order expansion theorem.
The mode-localization argument should be a separate theorem with
explicit hypotheses: uniform central-limit localization, a concave
principal profile, and a secondary-profile bound. Finally the
repository's logarithmic-core and remainder-transport organization can
be used for inversion, with its staircase caveat preserved
\cite{transseries}.
[Added October 1, 2026, batch 73O2: that is, the volume's
\texttt{p0:thm:lambert-core}, \texttt{p0:thm:perturbed-inversion} and
\texttt{p0:thm:staircase}~\cite{tai}; see the notes in
Section~\ref{tfp:sec:inverse}.] Existing repository module descriptions are a
roadmap, not evidence that these new hypotheses have already been
formalized.

\appendix
\section{A compact coefficient-generation specification}
The symbolic derivation in the package follows this finite recipe.
Generate $b_r(j)$ by \eqref{tfp:eq:cumrec} for $r\le10$ and sum their
polynomial parts. For even cumulants use nonnegative exponents
$(e_4,e_6,e_8,e_{10})$ satisfying
\begin{equation}
 e_4+2e_6+3e_8+4e_{10}\le4.
\end{equation}
Put $D=4e_4+6e_6+8e_8+10e_{10}$ and add
\begin{equation}
 \prod_{r\in\{4,6,8,10\}}
        \frac{\kappa_r^{e_r}}{e_r!(r!)^{e_r}}
   \frac{\He_D(d/\sqrt V)}{V^{D/2}}.
\end{equation}
For this precision, retaining the constant and quadratic terms of
each even Hermite polynomial suffices. Add the single odd term
$-\kappa_3d/(2V^2)$. Multiply by $e^{-d^2/(2V)}$, the variance factor
\eqref{tfp:eq:varfactor}, and the factorial factor \eqref{tfp:eq:stirling}.
Expand through $n^{-4}$, leaving $h_n,h_{2,n}$ as exact symbols.
Finally add $-18(-1)^{n+k}/n^4$.

The omitted odd terms begin at relative order $n^{-5}$:
$\kappa_5d/V^3=O(n^{-5})$ and
$\kappa_3^2/V^3=O(n^{-5})$ are the first examples.
The proper rational parts of $b_4(j)$ sum to $O(\log n)$ and first
affect this calculation at $O(\log n/n^6)$.
These order checks explain why polynomial cumulant parts suffice for
the four coefficients, although not for arbitrary order.

\section{Source audit and package contents}
The source check was performed on October 1, 2026. The live A039831
internal record displayed revision \#9, January 5, 2023, with the
conjecture still present. A039824 displayed its separate conjectures.
The ProveIt recursive-tree response inspected at the beginning of the
work identified commit
\texttt{a5ad4ac2e2d4e642a6c015836a70492d40fcbc71}.
Relevant repository reads were the root README and the consolidated
transseries/inversion README; targeted searches included the A039831
identifier and asymptotic inversion. Search indices and a live branch
need not share a commit, so the mathematical proofs above do not depend
on an unverified assertion about repository-wide coverage.

The accompanying archive contains the editable \LaTeX{} source, compiled
PDF, an exact coefficient and numerical verification program, a symbolic
coefficient derivation, an arbitrary-order finite-formula evaluator,
finite unit tests, numerical and inverse CSV
outputs, a source manifest, and build instructions. No external paper
or font file is redistributed. The bibliography supplies the primary
sources used for context and provenance.
[Added October 1, 2026, batch 73O2: in the shipped report the scripts are in
\texttt{code/}, their recorded outputs and the requirements file in
\texttt{data/}, the compiled PDF is a new build of this text, and the
delivered README is replaced by the report README;
Appendix~\ref{tfp:app:provenance} lists what is shipped.]

\section{Provenance and editorial notes}\label{tfp:app:provenance}
[Added October 1, 2026, batch 73O2.]

\paragraph{Source.}
The report was built from one manuscript, manuscript~62 of batch~73 of
ProveIt's incoming-reports intake: the archive
\texttt{A039831\_Two\_Fourier\_Peaks.zip} (main file \texttt{article.tex},
an 18-page PDF). It arrived in commit \texttt{aa43cc555} and was placed in
commit \texttt{6e193dd4f}, which deleted the archive from
\texttt{docs/incoming}; it survives in \texttt{aa43cc555}. Its pin is
commit \texttt{a5ad4ac2e} (the full hash is printed in the previous
section), an ancestor of the placement commit; the two repository files the
manuscript read, the root README and the transseries volume's README, are
unchanged between the pin and the placement.

\paragraph{What the write changed.}
\begin{itemize}
\item Every label gained the prefix \texttt{tfp:}; none was otherwise
 renamed or removed. The manuscript had 74 labels.
\item Added: the editorial note after the abstract; bracketed notes in
 Section~\ref{tfp:sec:inverse}, research question~9, Section~12 and
 Appendix~B; one bibliography entry~\cite{tai}; and this section.
\item Nothing was removed. Every non-claim of the manuscript is kept:
 A039824 is not proved; the remainders are asymptotic, not interval
 certificates; the numerical errors are measurements; the first-order mode
 prediction fails at $n=56$; the Dolgopyat--Hafouta theorem is not invoked;
 no global priority, referee review or formal verification is claimed.
\end{itemize}

\paragraph{A039824 is not resolved.}
The report proves Kote\v{s}ovec's conjecture on the \emph{largest}
coefficient (A039831). The \emph{number of distinct coefficient values}
(A039824), conjectured to be $n^2-3$ for $n>6$, is a different question, and
nothing here settles it.

\paragraph{Checks by the intake.}
The intake recomputed the rows exactly through $n=200$ ($M_1,\ldots,M_5=1,2,4,13,57$,
row masses $(n+1)!$) and found the relative error of the two-correction
formula \eqref{tfp:eq:leading} to be about $8.8\times10^{-4}$,
$1.24\times10^{-4}$ and $1.70\times10^{-5}$ at $n=50,100,200$, a ratio of
about $7.3$ per doubling, consistent with $O(\log n/n^3)$. It reran the three
programs on a copy; their outputs equal the shipped ones apart from line
endings and the elapsed-time line. These checks do not replace an
independent review of the proofs.

\begin{thebibliography}{9}
\bibitem{oeis831}
The OEIS Foundation, \emph{A039831: Largest coefficient in expansion of
$\prod_{i=1}^n(1+q+q^3+\cdots+q^{2i-1})$}.
Entry by Olivier G\'erard; asymptotic conjecture by V\'aclav Kote\v{s}ovec,
January 5, 2023. Inspected October 1, 2026.
\url{https://oeis.org/A039831}; internal record:
\url{https://oeis.org/A039831/internal}.

\bibitem{oeis824}
The OEIS Foundation, \emph{A039824: Number of different coefficient
values in the same product}. Conjectured eventual formula by Ralf
Stephan, March 7, 2004; related generating-function and recurrence
conjectures by Colin Barker, May 2, 2017. Inspected October 1, 2026.
\url{https://oeis.org/A039824}.

\bibitem{dolgopyat}
D. Dolgopyat and Y. Hafouta,
\emph{Edgeworth expansions for independent bounded integer valued
random variables}, Stochastic Processes and their Applications
\textbf{152} (2022), 486--532.
\url{https://arxiv.org/abs/2011.14852}.
The abstract and publication metadata were consulted for context;
the theorem is not invoked for our growing-support summands.

\bibitem{proveit}
V. Reshetnikov, \emph{ProveIt}, root README and repository tree,
inspected October 1, 2026.
\url{https://github.com/VladimirReshetnikov/ProveIt}.

\bibitem{transseries}
V. Reshetnikov, ProveIt research corpus,
\emph{Transseries: the polynomial--logarithmic calculus and its
inversions}, consolidated volume README, inspected October 1, 2026.
\url{https://github.com/VladimirReshetnikov/ProveIt/blob/main/Analysis/Transseries/docs/series-and-transseries/Transseries_And_Inversion/README.md}.

\bibitem{tai}\begingroup\raggedright
[Added in the batch-73O2 write.] V.~Reshetnikov and contributors,
\emph{Transseries and Inversion}, the volume whose README is~\cite{transseries}:
its source \path{transseries_and_inversion.tex} lies beside that README. Cited here: Part ``Inversion: the apparatus for a rapidly growing function''
(state of commit \texttt{6e193dd4f}).\par\endgroup
\end{thebibliography}
\end{document}
